% Part: proof-theory
% Chapter: natural-deduction
% Section: introduction

\documentclass[../../../include/open-logic-section]{subfiles}

\begin{document}

\olfileid{pt}{nat}{int}

\olsection{परिचय}

नैसर्गिक निगमन हा !!{proof}
पद्धतींचा एक परिवार आहे.
त्यात प्रत्येक तार्किक संयोजकासाठी
किंवा संख्यापकासाठी दोन
अनुमाने असतात: एक संकारकाचे
“प्रवेशन” करणारे
(!!a{introduction} नियम)
आणि दुसरे त्याचे “विलोपन”
करणारे (!!a{elimination} नियम).
उदाहरणार्थ, \Intro{\lif}
च्या निष्कर्षात
$!A \lif !B$ या रूपाची
!!{formula}s असतात,
तर \Elim{\lif} च्या
आधारविधानांत
$!A \lif !B$ या रूपाची
!!{formula}s असतात.

नैसर्गिक निगमनाच्या
पहिल्या दोन आवृत्त्या
१९३० च्या दशकात गेर्हार्ड
गेंट्सेन आणि स्तानिस्वाव
याश्कोव्स्की यांनी मांडल्या.
सिद्धता-उपपत्तीचे अभ्यासक
बहुधा गेंट्सेनच्या पद्धतीची
चर्चा करतात; अध्यापनात
सर्वाधिक वापरल्या जाणाऱ्या
पद्धतींचा उगम याश्कोव्स्कीच्या
पद्धतीत आहे. दोन्हीत
\emph{गृहीतकांपासून} सुरुवात
करता येते: अशी !!{formula}s
की ज्यांना !!{prove} करावे
लागत नाही, पण इतर
!!{formula}s !!{proving}
करताना वापरता येते.
संपूर्ण !!{proof} अशा
गृहीतकांवर अवलंबून असू शकते.
ही गृहीतके असिद्ध असल्यामुळे
ती !!{proof} तिचा निष्कर्ष
(तिचे \emph{अंतिम-!!{formula}})
सिद्ध करत नाही; निष्कर्ष
गृहीतकांपासून निष्पन्न होतो,
एवढेच ती दाखवते.

काही अनुमान नियमांचा
निष्कर्ष काही गृहीतकांवर
आता अवलंबून राहत नाही:
असे नियम गृहीतके “मुक्त”
करतात. उदाहरणार्थ,
\Intro{\lif} नियम
गृहीतक~$!A$ पासून
निष्कर्ष~$!B$ ची
!!a{proof} घेऊन
$!A \lif !B$ हा
निष्कर्ष देतो.
$!A \lif !B$ मध्ये
संपणारी !!{proof}
आता~$!A$ वर अवलंबून
नसते: ते गृहीतक
मुक्त झालेले असते.
गेंट्सेनच्या पद्धतीत
!!{proof}s हे
!!{formula}s चे वृक्ष
असतात; गृहीतके सर्वात
वरची !!{formula}s असतात,
आणि \Intro{\lif}
सारख्या नियमांवरील खुणा
कोणती गृहीतके मुक्त होतात
ते दाखवतात.
याश्कोव्स्कीच्या पद्धतीत
एखाद्या गृहीतकावर अवलंबून
असलेले !!{proof} चे भाग
वेगळे दाखवतात: उदाहरणार्थ,
उभी रेषा लावून आत सरकवलेले
किंवा चौकटीत बंद केलेले.
त्या चौकटीची पहिली ओळ
गृहीतक~$!A$ असते;
अंतिम ओळ सोपाधिक विधानाचे
उत्तरांग $!B$ असते;
$!A$ मुक्त केले जाते;
आणि \Intro{\lif} चा
निष्कर्ष~$!A \lif !B$
चौकटीबाहेर असतो.

या पद्धती “नैसर्गिक”
आहेत, कारण त्यांचे अनुमान
नियम आपण (निदान गणितज्ञ)
सहजपणे जसे युक्तिवाद करतो
त्याशी जुळतात. एखादे
सोपाधिक विधान सिद्ध
करण्यासाठी त्याचे पूर्वांग
गृहीत धरून त्यावरून उत्तरांग
सिद्ध करतो: तो \Intro{\lif}.
सोपाधिक विधान $!A \lif !B$
आणि त्याचे पूर्वांग~$!A$
माहीत असेल, तर~$!B$
निष्पन्न करतो: तो
\Elim{\lif}.
विकल्पयोग~$!A \lor !B$
गृहीत धरून काही निष्कर्ष
दाखवण्यासाठी प्रकरणांनुसार
सिद्धता वापरतो: तो
\Elim{\lor}.
सार्वत्रिक विधान
$\lforall[x][!A(x)]$
सिद्ध करण्यासाठी
कोणताही~$c$ घेऊन
$!A(c)$ सिद्ध करतो:
तो \Intro{\lforall}.
अशाच प्रकारे पुढे जाता येते.

उदाहरणार्थ, गेंट्सेन-पद्धतीच्या
नैसर्गिक निगमनात
$!A \lif (!A \lor !B)$
ची ही सोपी !!{proof} पाहा:
\[
\AxiomC{$\Discharge{!A}{x}$}
\RightLabel{\Intro{\lor}}
\UnaryInfC{$!A \lor !B$}
\DischargeRule{\Intro{\lif}}{x}
\UnaryInfC{$!A \lif (!A \lor !B)$}
\DisplayProof
\]
ती गृहीतक~$!A$ पासून
सुरू होते; \Intro{\lor}
वापरून~$!A \lor !B$
निष्पन्न करते; मग
\Intro{\lif} वापरून
$!A \lif (!A \lor !B)$
निष्पन्न करते.
\Intro{\lif} अनुमान
गृहीतक~$!A$ मुक्त करते;
खूण~$x$ ते दर्शवते.

सिद्धता-उपपत्तीचे अभ्यासक
!!{formula}s च्या वृक्षांवर
काम करणाऱ्या गेंट्सेनच्या
मूळ पद्धती पसंत करतात;
तरी एक रूपांतरही
सिद्धता-उपपत्तीत महत्त्वाचे आहे.
गृहीतके~$\Gamma$ यांपासून
$!A$ ची !!^a{proof}
म्हणजे $!A$ हे~$\Gamma$
पासून निष्पन्न होते,
अर्थात $\Gamma \Entails !A$,
हे दाखवते.
नैसर्गिक निगमनाचे अनुमान
नियम अशा निष्पन्नतेबद्दल
विधान करतात असेही
स्वाभाविकपणे वाचता येते.
उदाहरणार्थ, \Intro{\lif}
म्हणतो की जर $\Gamma + !A
\Entails !B$, तर
$\Gamma \Entails !A \lif !B$.
म्हणून नियम थेट गृहीतके
आणि त्यांपासून निष्पन्न
होणारा निष्कर्ष या दोहोंवर
लागू होतील असे त्यांचे
रूप देता येते: म्हणजे
ते $\Gamma \Sequent !A$
या क्रमवर्तींवर काम करतात.
वरील सोपी सिद्धता अशा
पद्धतीत पुढीलप्रमाणे दिसेल:
\[
\Axiom$x:!A \fCenter !A$
\RightLabel{\Intro{\lor}}
\UnaryInf$x:!A \fCenter !A \lor !B$
\DischargeRule{\Intro{\lif}}{x}
\UnaryInf$\fCenter !A \lif (!A \lor !B)$
\DisplayProof
\]
येथे नियम तेच आहेत:
ते~$\Sequent$ च्या
उजवीकडील !!{formula}
वर काम करतात;
डावीकडील !!{formula}s
प्रत्येक टप्प्यावर कोणत्या
गृहीतकांवर अवलंबित्व आहे
ते फक्त दाखवतात.

!!{introduction}/!!{elimination}
या नियम-जोड्यांचा पूर्ण संच
स्वतःहून अभिजात तर्कशास्त्रासाठी
पूर्ण नसतो. त्यात~$\lfalse$
शी संबंधित आणखी दोन
नियम जोडावे लागतात.
पहिला “अंतःप्रज्ञावादी
असत्यता नियम”~\FalseInt,
म्हणजे “स्फोट”:
$\lfalse$ पासून काहीही
निष्पन्न करता येते.
दुसरा अप्रत्यक्ष सिद्धतेचा
अभिजात सिद्धांत~\FalseCl{}
(किंवा “अभिजात असत्यता नियम”):
$\lnot !A$ पासून
$\lfalse$ !!{prove}
करता येत असेल, तर
$!A$ !!{prove}d झाले.
\FalseCl वगळल्यास
अंतःप्रज्ञावादी तर्कशास्त्राला
योग्य अशी पद्धत मिळते.
दोन्ही नियम वगळल्यास
“किमान तर्कशास्त्र”
मिळते.

अभिजात आणि
अंतःप्रज्ञावादी तर्कशास्त्रासाठी
गेंट्सेन-पद्धतीच्या
नैसर्गिक निगमनाची आपण
चर्चा करू. या
\Log{N1c} आणि \Log{N1i}
पद्धती आहेत.
\Log{N2c} आणि \Log{N2i}
या संबंधित पद्धती
केवळ !!{formula}s~$!A$
यांवर काम करण्याऐवजी
$\Gamma \Sequent !A$
या क्रमवर्तींवर काम करतात.

\end{document}
